\documentclass[a4paper,11pt]{article}
\usepackage[a4paper,margin=2cm]{geometry}
\usepackage[T1]{fontenc}
\usepackage[utf8]{inputenc}
\usepackage[ngerman,english]{babel}
\usepackage{lmodern,microtype,xcolor,parskip,enumitem}
\usepackage[most]{tcolorbox}
\usepackage{titlesec}
\usepackage[hidelinks]{hyperref}
\definecolor{HeaderBlueDark}{RGB}{70,110,170}
\definecolor{RowBlueLight}{RGB}{235,243,255}
\definecolor{RowBlueDark}{RGB}{220,233,250}
\definecolor{SoftGray}{RGB}{90,90,90}
\setlength{\parindent}{0pt}
\setlist[itemize]{leftmargin=1.4em,itemsep=0.35em,topsep=0.35em}
\linespread{1.06}
\titleformat{\section}[block]{\Large\bfseries\color{HeaderBlueDark}}{}{0pt}{}
\titlespacing{\section}{0pt}{1.3em}{0.8em}
\newtcolorbox{exbox}{enhanced,breakable,colback=RowBlueLight,colframe=RowBlueDark,boxrule=0.4pt,arc=1.5mm,left=1.2mm,right=1.2mm,top=1mm,bottom=1mm,title={Beispiele (Examples)},fonttitle=\bfseries,colbacktitle=RowBlueDark,coltitle=black}
\NewDocumentEnvironment{grammarentry}{mm+b}{\begin{tcolorbox}[enhanced,breakable,colback=white,colframe=HeaderBlueDark,boxrule=0.6pt,arc=1.5mm,left=1.5mm,right=1.5mm,top=1mm,bottom=1mm,colbacktitle=HeaderBlueDark,coltitle=white,fonttitle=\bfseries,title={#1\hfill\normalfont\footnotesize #2}]#3\end{tcolorbox}}{}
\newcommand{\GermanText}[1]{\textbf{Deutsch: }#1\par}
\newcommand{\EnglishText}[1]{{\small\color{SoftGray}\textbf{English: }#1\par}}
\newcommand{\ExampleItem}[2]{\item \textbf{#1}\\{\small\color{SoftGray}#2}}
\title{Das Fragewort „wie“ (The question word “wie”)}
\date{}
\begin{document}
\maketitle
\tableofcontents
\newpage
\section{Grundbedeutung und Wortart (Basic meaning and word class)}
\begin{grammarentry}{wie}{Interrogativadverb (modal)}
\GermanText{\textbf{wie} ist ein modales Interrogativadverb. Es fragt nach Art und Weise, Zustand oder Beschaffenheit: „auf welche Weise?“ oder „was für ...?“.}
\EnglishText{\textbf{wie} is a modal interrogative adverb. It asks about manner, state, or quality: “in what way?” or “what ... like?”.}

\end{grammarentry}
\section{Satzstellung und Verwendung (Word order and use)}
\begin{grammarentry}{Direkte und indirekte Frage}{Fragewortstellung}
\GermanText{In einer direkten W-Frage steht das finite Verb gewöhnlich an zweiter Stelle nach dem Fragewort. In einer indirekten Frage steht es am Ende. Mit einem Adjektiv fragt \textbf{wie} nach einem Merkmal oder Grad. In indirekten Fragen steht das finite Verb am Ende.}
\EnglishText{In a direct information question, the finite verb normally follows the question word in second position. In an indirect question, it goes at the end. With an adjective, \textbf{wie} asks about a quality or degree. In indirect questions, the finite verb goes at the end.}
\end{grammarentry}
\begin{exbox}\begin{itemize}
\ExampleItem{Wie funktioniert diese Maschine?}{How does this machine work?}
\ExampleItem{Wie alt ist dein Bruder?}{How old is your brother?}
\ExampleItem{Kannst du mir sagen, wie das Gerät funktioniert?}{Can you tell me how the device works?}
\end{itemize}\end{exbox}
\section{Zusammenfassung (Summary)}
\begin{grammarentry}{Merksatz}{kurz und wichtig}
\GermanText{Merksatz: \textbf{wie} fragt nach Art, Zustand oder Beschaffenheit.}
\EnglishText{Summary: \textbf{wie} asks about manner, state, or quality.}
\end{grammarentry}
\end{document}
